\begin{pastexam}[label={px:df3-108-ds-04}]{地方三等資訊處理 108 年 第四題}
（申論；20 分）

下面的無向圖（undirected graph）表示四個人的關係，如張三與李四有關係，這二人之間有邊（edge）相連，則可走訪。
括弧內為人名縮寫，如張三（Chang San）的縮寫為 CS。若同時有兩個以上的人可處理，則先處理人名縮寫的字母順序
較小者。
\begin{center}
\begin{tikzpicture}[v/.style={font=\small,align=center}]
\node[v] (LS) at (0,1.1) {李四(LS)};
\node[v] (CS) at (-3,0) {張三\\(CS)};
\node[v] (WW) at (3,0) {王五\\(WW)};
\node[v] (CL) at (0,-1.1) {趙六(CL)};
\draw (CS) -- (LS) (LS) -- (WW) (LS) -- (CL) (CS) -- (CL) (CL) -- (WW);
\end{tikzpicture}
\end{center}
\begin{itemize}[leftmargin=2.2em]
\item[㈠] 由張三（CS）出發，用佇列（queue）做廣度優先搜尋（breadth-first search）走訪所有人，請寫出走訪順序的
中文人名。（10 分）
\item[㈡] 由張三（CS）出發，用堆疊（stack）做深度優先搜尋（depth-first search）走訪所有人，請寫出走訪順序的
中文人名。（10 分）
\end{itemize}
\end{pastexam}

\begin{pastsoln}
答案：㈠ 張三、趙六、李四、王五；㈡ 張三、趙六、李四、王五。

縮寫的字母順序是 $\mathrm{CL} < \mathrm{CS} < \mathrm{LS} < \mathrm{WW}$。 鄰接串列（依字母排序）：
\begin{center}
\begin{tabular}{ll}
\toprule
人 & 相鄰的人 \\
\midrule
張三 CS & CL、LS \\
趙六 CL & CS、LS、WW \\
李四 LS & CL、CS、WW \\
王五 WW & CL、LS \\
\bottomrule
\end{tabular}
\end{center}

\textbf{㈠ 廣度優先}（資結書第~21.3~節）。 放入佇列時就標記，每次取出佇列前端的人，把還沒標記的鄰居依字母順序
放進佇列：
\begin{center}
\begin{tabular}{lll}
\toprule
取出走訪 & 放入 & 佇列（左端是前端） \\
\midrule
& CS & CS \\
CS（張三） & CL、LS & CL, LS \\
CL（趙六） & WW（CS、LS 已標記） & LS, WW \\
LS（李四） & 無 & WW \\
WW（王五） & 無 & 空 \\
\bottomrule
\end{tabular}
\end{center}
走訪順序：張三、趙六、李四、王五。

\textbf{㈡ 深度優先}（資結書第~21.4~節、虛擬碼~21.4）。 每次 pop 一個人，沒走過就走訪，再把他還沒走過的鄰居
\textbf{反著字母順序} push，讓字母小的在頂端、先被處理（這樣才符合「先處理字母順序較小者」）：
\begin{center}
\begin{tabular}{lll}
\toprule
pop & 走訪 & push 之後的堆疊（右端是頂端） \\
\midrule
& & CS \\
CS & 張三 & LS, CL \\
CL & 趙六 & LS, WW, LS \\
LS & 李四 & LS, WW, WW \\
WW & 王五 & LS, WW \\
WW、LS & （都走過了，丟掉） & 空 \\
\bottomrule
\end{tabular}
\end{center}
走訪順序：張三、趙六、李四、王五。 用遞迴的想法驗證：從張三出發，往字母最小的鄰居趙六走；趙六還沒走過的鄰居
是李四、王五，先走李四；李四還沒走過的只剩王五，走到王五，四人走完。

這張圖兩種走法的順序剛好相同，但過程不同：BFS 是「先把張三的鄰居趙六、李四都排進來」，李四是因為和張三相鄰
才排第三；DFS 是沿著 張三 $\to$ 趙六 $\to$ 李四 $\to$ 王五 一路走下去，李四是從趙六走過去的。 兩者的走訪樹不同：
BFS 樹的邊是 CS–CL、CS–LS、CL–WW；DFS 樹的邊是 CS–CL、CL–LS、LS–WW。
\end{pastsoln}
